\section{Conclusiones}
\label{cap:conclusiones}

El proyecto construyó una arquitectura multimodelo completa sobre un volcado real de ListenBrainz de 4\,855\,809 escuchas: MongoDB con los documentos originales, Cassandra con la serie temporal y los agregados, Neo4j con el grafo de relaciones y un DAG de Airflow que mantiene las tres bases a partir de la misma fuente y calcula cinco indicadores. Todo se levanta con un solo \code{docker-compose.yml}.

\subsection{Aprendizajes}

Aprendimos que el modelo de datos se diseña a partir de las consultas. En Cassandra eso es explícito: una tabla por consulta, particiones por usuario y mes para que ninguna crezca sin límite y el ranking guardado como clustering porque es la única forma de que Cassandra lo devuelva ordenado. A cambio, las consultas de Cassandra tardan entre 6 y 16~ms sobre millones de filas. En Neo4j el mismo criterio llevó a agregar las escuchas en relaciones con cantidad y proporción, que redujeron 4,86 millones de eventos a 1,79 millones de relaciones usuario--canción y dejaron en el grafo exactamente lo que necesitan las recomendaciones.

En un pipeline que se reintenta, la idempotencia importa más que la velocidad. Las tres cargas se pueden repetir sin duplicar datos por razones distintas: huellas SHA-256 y checkpoint en MongoDB, ids deterministas y registro de lotes en Cassandra, \texttt{MERGE} y estado del lote en Neo4j. En Cassandra esto incluso obligó a abandonar los counters, que suman dos veces si un lote se repite, y a recalcular los agregados en cada carga.

También vimos que hay que conocer el dato antes de interpretarlo. Un incremental describe la ingesta, no la escucha: hay escuchas fechadas en 2005 enviadas en 2026, el 26 de octubre concentra el 28~\% del archivo y el volcado se corta el 27 a las 00:01. Sin esa lectura, el crecimiento de $+3\,364$~\% o la participación de BTS parecerían hallazgos sobre los oyentes, cuando son efectos de cómo llegaron los datos.

\subsection{Trabajo futuro}

Como siguiente paso, \code{extraer_datos} puede descargar el siguiente incremental publicado por ListenBrainz, Neo4j puede cargarse de forma incremental, Cassandra puede recalcular solo los días y meses que toca cada lote, y los artistas pueden normalizarse con los identificadores de MusicBrainz.

\subsection{Repositorio y video}

El código fuente completo (loaders, scripts de creación de las bases, DAG de Airflow y dashboard) está en el repositorio \repolink. El video demostrativo (carpeta de Google Drive) está disponible en \videolink.
